\documentclass[
    language=chinese,
    doctype=thesis,
    institution=none,
    titlestyle=thesis,
]{../../omnilatex}

\RequirePackage{config/document-settings}

\begin{document}

\frontmatter
\maketitle
\tableofcontents

\mainmatter
\chapter{引言}
本研究旨在探讨中文语境下的学术写作与排版规范。随着中文学术出版的国际化，高质量的中文排版需求日益增长。本文介绍基于 OmniLaTeX 的中文论文模板，涵盖字体配置、章节结构和参考文献管理。

\chapter{研究背景}
中文排版具有独特的规范要求：段落首行缩进两个字符宽度，行距一般为字号的 1.5 至 2 倍，标点符号使用全角形式。OmniLaTeX 基于 LuaLaTeX 和 luatexja 引擎，自动处理这些细节，无需手动调整。

在当今全球化的学术环境中，中文技术文档的质量直接影响到学术交流的效率。一份排版规范、格式统一的论文，不仅降低了读者的认知负担，也为跨国学术合作搭建了桥梁。

\chapter{方法论}
本研究采用文献分析与实证研究相结合的方法。通过对比分析国内外中文排版规范，总结出适用于学术论文的排版准则。实验数据表明，使用 OmniLaTeX 排版的中文论文在可读性和专业性方面均优于传统排版工具。

实验在以下方面进行了对比：字体渲染质量、中英混排间距、标点压缩处理以及参考文献格式化。结果表明 OmniLaTeX 在所有维度上均达到或超过了现有商业排版软件的水准。

\chapter{实验结果}
实验结果显示，OmniLaTeX 的中英混排间距处理机制显著提升了文档的视觉效果。在 $O(n \log n)$ 时间复杂度的算法描述中，中英文字符之间的自动间距调整确保了公式的清晰可读。

欧拉公式 $e^{i\pi} + 1 = 0$ 和高斯积分 $\int_{-\infty}^{+\infty} e^{-x^2} \, dx = \sqrt{\pi}$ 等数学表达在中文语境下的排版效果良好，符号间距均匀，基线对齐准确。

\chapter{结论}
本文提出了基于 OmniLaTeX 的中文论文排版方案。通过 LuaLaTeX 和 luatexja 的深度集成，实现了简体中文、繁体中文以及中英混排内容的原生支持。该方案可广泛应用于学位论文、学术期刊和技术报告的排版。

未来工作将聚焦于多栏排版支持、中文摘要自动生成以及与主流文献管理工具的集成。

\backmatter
\printbibliography

\end{document}
